\section{Proof details and auxiliary checks}
\label{app:proofs}

\subsection{Orbit indicators}

If a permutation $p$ has cycles of length dividing $d$ and $d$ is invertible
in the coefficient field, then
\[
  I+p+\cdots+p^{d-1}
\]
sends every basis vector in a cycle of exact length $e$ to $(d/e)$ times the
cycle indicator.  Hence its image is exactly the span of orbit indicators.
This proves the relation-image statement used in \cref{thm:rank} for
$(S,d)=(S,2)$ and $(R,d)=(R,3)$.

\subsection{Connectedness of the incidence dessin}

Two state edges in $D_n$ lie in the same connected component precisely when
one can pass between their states by a word in $S$ and $R$.  The projective
modular action generated by $S[a:b]=[-b:a]$ and $R[a:b]=[-b:a+b]$ is
transitive on $P^1(\ZZ/n\ZZ)$.  Therefore $D_n$ is connected.

\subsection{Finite-support versus unrestricted cohomology}

The global complex contains infinitely many blocks.  Cellular chains are
finite sums by definition, so the homological primitive ledger is
\[
  H_1=\bigoplus_{n\ge2}M_n^* .
\]
If one instead takes unrestricted global cochains, then
\[
  H^1=\prod_{n\ge2}M_n .
\]
The project uses the finite-support cohomology ledger
$L_{\rm fs}=\bigoplus_{n\ge2}M_n$ only as a ledger, not as ordinary global
cohomology and not as determinant ownership.

\subsection{Rank field}

The rank audit uses $\mathbb F_{1000003}$ only to validate sparse linear
algebra at finite cutoffs.  The theorems are characteristic-zero statements.
The chosen prime avoids the characteristics $2$ and $3$ in which the displayed
relation images would degenerate.
